\documentclass{article}
\usepackage[top=20mm,left=20mm,right=15mm]{geometry}
\begin{document}
\section{Conclusion}

\subsection{Summary of Findings}

This project successfully designed and built a prototype of an automatic streetlight control and fault detection system. 
The system uses an ESP32-WROOM-32E microcontroller connected to an LDR sensor for ambient light detection, an ACS712 current sensor for fault detection, a relay module for switching the streetlight, 
and a GSM module for sending SMS alerts.
All components were powered from the mains supply through an HLK-20M05 power module and an AMS1117-3.3 voltage regulator. 
The complete circuit was drawn in a schematic editor and implemented on a custom two-layer PCB

Testing confirmed that the system correctly switches the streetlight on at night and off during the day, and that it can detect bulb faults and send an SMS alert within a few seconds. The total component cost was kept below 7,000 LKR, 
making the system affordable for small local government bodies.

\subsection{Implications}

This project shows that it is possible to build an effective, affordable, and automated streetlight management system using low-cost, widely available components.
If deployed at scale, such a system could significantly reduce the energy wasted on daytime streetlight operation and reduce the delay between fault occurrence and fault repair. 
Both of these outcomes would have a positive impact on public safety and on the budgets of local government bodies.

\subsection{Future Work}

The following improvements are planned for future versions of the system:

\begin{itemize}
    \item \textbf{Wi-Fi / Cloud Integration:} The ESP32 already has built-in Wi-Fi capability. A future version of the firmware could send fault alerts and operational data to a cloud-based dashboard, allowing a single operator to monitor all streetlights in an area from one screen.
    \item \textbf{Dimming Control:} Instead of simply switching the light on or off, the system could be extended to dim the light during low-traffic hours using pulse-width modulation (PWM), further reducing energy consumption.
    \item \textbf{Solar Power:} For areas where mains power is unreliable, the system could be adapted to run from a solar panel and battery.
    \item \textbf{Multiple Fault Types:} The current system detects open-circuit faults (bulb failure). Future versions could also detect short-circuit faults and over-current conditions by adding more sophisticated signal processing to the ACS712 readings.
    \item \textbf{GPS Location Reporting:} Adding a GPS module would allow the SMS alert to include the location of the faulty streetlight, making it easier for maintenance teams to find and fix the problem quickly.
\end{itemize}
\end{document}